On prépare une solution aqueuse d'acide azélaïque en dissolvant la masse
$m=\qty{0.188}{\g}$ de cet acide dans le volume $V=\qty{100}{\mL}$ d'eau
distillée. La mesure de $\mathrm{pH}$ de la solution à
$\qty{25}{\degreeCelsius}$ donne $\mathrm{pH}=3,28$.

\begin{donnees}
    \begin{itemize}
        \item $M(\ce{AH})=\qty{188}{\g\per\mol}$
    \end{itemize}
\end{donnees}

\begin{enumerate}
    \item Écrire l'équation qui modélise la réaction entre l'acide $\ce{AH}$ et
          l'eau.
    \item Dresser le tableau d'avancement de la réaction entre l'acide azélaïque
          et l'eau.
    \item Montrer que le taux d'avancement final de la réaction s'écrit
          $\tau=\frac{M\cdot V\cdot 10^{-\mathrm{pH}}}{m}$.
    \item Calculer la valeur de $\tau$. Déduire.
    \item Montrer que l'expression du quotient de la réaction $Q_{r,\eq}$ à
          l'équilibre s'écrit
          $Q_{r,\eq}=\frac{m\cdot\tau^{2}}{M\cdot V(1-\tau)}$.
    \item Vérifier que la valeur du $\mathrm{pK_{A}}$ du couple
          $(\ce{AH_{(aq)}}\ttslash{}\ce{A^-_{(aq)}})$ est
          $\mathrm{pK_{A}}=4,54$.
    \item Déterminer, parmi les espèces $\ce{AH_{(aq)}}$ et $\ce{A^-_{(aq)}}$,
          l'espèce qui prédomine dans la solution à l'équilibre.
\end{enumerate}

